\begin{figure*}

  \begin{panel}{(a)}{\textwidth}
    \hspace{-2mm}%
    \setlength\plotwidth{\textwidth}%
    \setlength\plotheight{36mm}%
    \tikzsetnextfilename{qpaint_trace_n10}%
\begin{tikzpicture}
  \argsinput[
    histogramcsv=figures/data/qpaint_kinetics/intensity_histogram.csv,
    tracecsv=figures/data/qpaint_kinetics/trace_and_fit_N10.csv,
    traceintensitycol=trace,
    zcol=fit,
    histbincol=N10_bins,
    histcountcol=N10_measured,
    modelfitcol=N10_model,
    plotlabel={$\noexpand\n=10$}
  ]{figures/plots/fancy_trace.plot}
  \setlength\plotwidth{0.1\plotwidth}
  \setlength\plotheight{\plotwidth}
  \pgfplotsset{every axis/.style={
    anchor=below south east,
    at={($(histogram.south east)-(2mm,-2mm)$)},
    xtick distance=2
  }}
  \argsinput[
    posteriorcsv=figures/data/qpaint_kinetics/posteriors.csv,
    posteriorcol=posterior_10,
    noxlabel
  ]{figures/plots/p_n_given_x.plot}
  \node[above=0mm of posterior] {posterior};
\end{tikzpicture}

    \vspace{-4mm}
  \end{panel}

  \begin{panel}{(b)}{\textwidth}
    \hspace{-2mm}%
    \setlength\plotwidth{\textwidth}%
    \setlength\plotheight{36mm}%
    \tikzsetnextfilename{qpaint_trace_n20}%
\begin{tikzpicture}
  \argsinput[
    histogramcsv=figures/data/qpaint_kinetics/intensity_histogram.csv,
    tracecsv=figures/data/qpaint_kinetics/trace_and_fit_N20.csv,
    traceintensitycol=trace,
    zcol=fit,
    histbincol=N20_bins,
    histcountcol=N20_measured,
    modelfitcol=N20_model,
    plotlabel={$\noexpand\n=20$}
  ]{figures/plots/fancy_trace.plot}
  \setlength\plotwidth{0.1\plotwidth}
  \setlength\plotheight{\plotwidth}
  \pgfplotsset{every axis/.style={
    anchor=below south east,
    at={($(histogram.south east)-(2mm,-2mm)$)},
    xtick distance=2
  }}
  \argsinput[
    posteriorcsv=figures/data/qpaint_kinetics/posteriors.csv,
    posteriorcol=posterior_20,
    noxlabel
  ]{figures/plots/p_n_given_x.plot}
  \node[above=0mm of posterior] {posterior};
\end{tikzpicture}

    \vspace{-4mm}
  \end{panel}

  \begin{panel}{(c)}{\textwidth}
    \setlength\plotwidth{72mm}%
    \setlength\plotheight{72mm}%
    \centering%
    \input{figures/panels/qpaint_counting_estimates.tikz}
  \end{panel}

  \caption{
    \panelref{a, b} Simulated traces and \ours fits, \pon=0.006, \poff=0.2,  and  $n$ = 10/20 emitters. 
    \ours was fit with strong priors on all parameters. Even though we are in an unfavorable regime, and 
    the histograms of these traces have faw fewer peaks than the true count, \ours is still able to count accurately.
    %
    \panelref{c} Above $n > 15$ \qpaint underestimates the count, while 
    the posterior of \ours remains accurate on average. \textit{colorbar} shows the 
    average probability across 10 traces.
  }
  \label{fig:results:qpaint_counting}
\end{figure*}

